% Caption for cand_c19_all_tasks_calibrated.png.

\caption{\textbf{Success against calibrated actuator energy, and where the
schedules sit on the plane.} Left: every operating point of all four workloads,
energy converted to joules per episode and plotted on a log axis. Calibration is
global and applied per episode, $E = c\,\int\!\sum\tau_{\mathrm{grav}}^2\,dt +
P_{\mathrm{idle}}T$, reduced as elsewhere (median over an arm's 24 episodes, mean
over seeds). For widowx, $c=0.827\,\mathrm{W/(N\,m)^2}$ and
$P_{\mathrm{idle}}=4.61$\,W are derived from ROBOTIS' published stall torque and
current for the Dynamixel XM430-W350. \emph{The two google series are hollow
because their constants are a proxy}---Everyday Robots published no actuator
data, $P_{\mathrm{idle}}$ accounts for roughly 92\% of their result, and
calibration does not preserve their ordering (Spearman between the raw $\tau^2$
channel and calibrated energy is $+0.26$ on coke and $-0.04$ on drawer, against
$+0.92$ and $+0.77$ on widowx); they are indicative only. Stars mark each
workload's best schedule by success rate. Right: success rate on the
(period, window) grid, one small multiple per workload, each on its own
gradient; hatched cells are propagated from a twin with an identical schedule,
$\times$ and ? mark infeasible and unsolved within the CP-SAT budget. Plotted
energy is the quasi-static lower bound; the raw $\tau^2$ proxy cannot be
converted at face value because the simulated PD torque runs about 76$\times$ the
physical requirement. Derivation in \texttt{calib/ENERGY\_CALIBRATION.md}.}
